% Собран калькулятором Weconomics. Компилируется обычным pdflatex.
% Модель: Неравенство доходов.
\documentclass[12pt,a4paper]{article}
\usepackage[T2A]{fontenc}
\usepackage[utf8]{inputenc}
\usepackage[english,russian]{babel}
\usepackage{amsmath}
\usepackage{pgfplots}
\pgfplotsset{compat=1.18}
\usetikzlibrary{arrows.meta}
\usepackage{geometry}
\geometry{margin=2cm}
\definecolor{c2D2D2D}{HTML}{2D2D2D}
\definecolor{c6B6457}{HTML}{6B6457}
\definecolor{c5C5346}{HTML}{5C5346}
\definecolor{c8C8C84}{HTML}{8C8C84}
\definecolor{c2F6FED}{HTML}{2F6FED}
\definecolor{c2D6EED}{HTML}{2D6EED}
\definecolor{c75756D}{HTML}{75756D}
\definecolor{c2E9E44}{HTML}{2E9E44}
\definecolor{c27863A}{HTML}{27863A}
\definecolor{cE0563B}{HTML}{E0563B}
\definecolor{cD13E21}{HTML}{D13E21}
\begin{document}
\begin{figure}[h]
\centering
\begin{tikzpicture}
% -- Панель lorenz: x от 0 до 100, y от 0 до 100
\begin{axis}[name=panelA, at={(46.5pt,33.0pt)}, anchor=south west, width=226.5pt, height=226.5pt, scale only axis, clip=false, xmin=0, xmax=100, ymin=0, ymax=100, axis x line=middle, axis y line=middle, axis line style={c2D2D2D, line width=0.68pt, -{Stealth[length=5pt]}}, tick align=outside, major tick length=3pt, scaled ticks=false, tick style={c2D2D2D, line width=0.4pt}, xtick={10,20,30,40,50,60,70,80,90,100}, ytick={20,40,60,80}, tick label style={font=\fontsize{8}{9.2}\selectfont, text=c6B6457}, extra y ticks={100}, extra y tick labels={}, extra tick style={grid=major, major tick length=0pt}, grid=major, grid style={line width=0.40pt, c5C5346, opacity=0.13}]
\node[anchor=base, inner sep=0pt, font=\fontsize{9}{10.3}\selectfont, text=c6B6457, yshift=-28.5pt] at (axis cs:50,0) {$N, \%$};
\node[anchor=base west, inner sep=0pt, font=\fontsize{9}{10.3}\selectfont, text=c6B6457, yshift=3.8pt] at (axis cs:1.987,100) {$I, \%$};
\begin{scope}
\clip (axis cs:0,0) rectangle (axis cs:100,100);
\fill[c8C8C84, opacity=0.15] (axis cs:0,0) -- (axis cs:20,8) -- (axis cs:40,20) -- (axis cs:60,36) -- (axis cs:80,60) -- (axis cs:100,100) -- (axis cs:100,0) -- cycle;   % область: многоугольник, вершины точные
\fill[c2F6FED, opacity=0.12] (axis cs:0,0) -- (axis cs:20,8) -- (axis cs:40,20) -- (axis cs:60,36) -- (axis cs:80,60) -- (axis cs:100,100) -- cycle;   % область: многоугольник, вершины точные
\node[anchor=base west, inner sep=0pt, font=\fontsize{10}{11.5}\selectfont, text=c2D6EED] at (axis cs:33,52) {$A$};
\node[anchor=base west, inner sep=0pt, font=\fontsize{10}{11.5}\selectfont, text=c75756D] at (axis cs:66,18) {$B$};
\draw[c2E9E44, line width=0.81pt, dash pattern=on 3.6pt off 2.4pt] (axis cs:0,0) -- (axis cs:100,100);
\node[anchor=base west, inner sep=1pt, font=\fontsize{9}{10.3}\selectfont, text=c27863A] at (axis cs:80,86) {Равенство};
\addplot[c2F6FED, line width=1.17pt, forget plot] coordinates {(0,0) (20,8) (40,20) (60,36) (80,60) (100,100)};   % БЕЗ ЗАПИСИ: ломаная, вершины точные
\draw[cE0563B, line width=1.17pt] (axis cs:60,36) -- (axis cs:60,60);
\fill[cE0563B] (axis cs:60,36) circle[radius=1.8pt];
\fill[cE0563B] (axis cs:60,60) circle[radius=1.8pt];
\node[anchor=base west, inner sep=1pt, font=\fontsize{9}{10.3}\selectfont, text=cD13E21] at (axis cs:62.318,48) {Робин Гуд $0{,}24$};
\end{scope}
\fill[c2D2D2D] (axis cs:20,8) circle[radius=2.4pt];
\fill[c2D2D2D] (axis cs:40,20) circle[radius=2.4pt];
\fill[c2D2D2D] (axis cs:60,36) circle[radius=2.4pt];
\fill[c2D2D2D] (axis cs:80,60) circle[radius=2.4pt];
\node[anchor=north east, inner sep=0pt, xshift=-4pt, yshift=-4pt, font=\fontsize{8}{9.2}\selectfont, text=c6B6457] at (axis cs:0,0) {$0$};
\end{axis}
\end{tikzpicture}
\end{figure}
\end{document}